\title{2021 年高考数学 北京卷}
\date{}
\maketitle

\section*{一、选择题}
本题共10小题，每小题4分，共40分。在每小题列出的四个选项中，选出符合题目要求的一项。

\begin{question}[points=4]
已知集合 $A=\{x\mid -1<x<1\}$, $B=\{x\mid 0\leq x\leq 2\}$, 则 $A\cup B=$ （ ）
\begin{choices}
  \item $(-1,2)$
  \item $(-1,2]$
  \item $[0,1)$
  \item $[0,1]$
\end{choices}
\begin{solution}
\textbf{答案：}B

\textbf{分析：}结合题意利用并集的定义计算即可.

\textbf{详解：}由题意可得 $A\cup B=\{x\mid -1<x\leq 2\}$, 即 $A\cup B=(-1,2]$.
\end{solution}
\end{question}

\begin{question}[points=4]
在复平面内, 复数 $z$ 满足 $(1-i)z=2$, 则 $z=$ （ ）
\begin{choices}
  \item $2+i$
  \item $2-i$
  \item $1-i$
  \item $1+i$
\end{choices}
\begin{solution}
\textbf{答案：}D

\textbf{分析：}由题意利用复数的运算法则整理计算即可求得最终结果.

\textbf{详解：}由题意可得 $z=\frac{2}{1-i}=\frac{2(1+i)}{(1-i)(1+i)}=\frac{2(1+i)}{2}=1+i$.
\end{solution}
\end{question}

\begin{question}[points=4]
已知 $f(x)$ 是定义在 $[0,1]$ 上的函数, 那么“函数 $f(x)$ 在 $[0,1]$ 上单调递增”是“函数 $f(x)$ 在 $[0,1]$ 上的最大值为 $f(1)$”的（ ）
\begin{choices}
  \item 充分而不必要条件
  \item 必要而不充分条件
  \item 充分必要条件
  \item 既不充分也不必要条件
\end{choices}
\begin{solution}
\textbf{答案：}A

\textbf{分析：}利用两者之间的推出关系可判断两者之间的条件关系.

\textbf{详解：}若函数 $f(x)$ 在 $[0,1]$ 上单调递增, 则 $f(x)$ 在 $[0,1]$ 上的最大值为 $f(1)$; 若 $f(x)$ 在 $[0,1]$ 上的最大值为 $f(1)$, 比如 $f(x)=\left(x-\frac{1}{3}\right)^2$, 但 $f(x)=\left(x-\frac{1}{3}\right)^2$ 在 $\left[0,\frac{1}{3}\right]$ 为减函数, 在 $\left[\frac{1}{3},1\right]$ 为增函数, 故 $f(x)$ 在 $[0,1]$ 上的最大值为 $f(1)$ 推不出 $f(x)$ 在 $[0,1]$ 上单调递增, 故“函数 $f(x)$ 在 $[0,1]$ 上单调递增”是“$f(x)$ 在 $[0,1]$ 上的最大值为 $f(1)$”的充分不必要条件.
\end{solution}
\end{question}

\begin{question}[points=4]
某四面体的三视图如图所示, 该四面体的表面积为（ ）

\begin{center}
  \includegraphics[width=.38\linewidth]{question-4-1.jpg}
\end{center}
\begin{choices}
  \item $\frac{3+\sqrt{3}}{2}$
  \item $4$
  \item $3+\sqrt{3}$
  \item $2$
\end{choices}
\begin{solution}
\textbf{答案：}A

\textbf{分析：}根据三视图可得如图所示的几何体（三棱锥）, 根据三视图中的数据可计算该几何体的表面积.

\textbf{详解：}根据三视图可得如图所示的几何体-正三棱锥 $O-ABC$, 其侧面为等腰直角三角形, 底面等边三角形, 由三视图可得该正三棱锥的侧棱长为 $1$, 故其表面积为 $3\times\frac{1}{2}\times1\times1+\frac{\sqrt{3}}{4}\times(\sqrt{2})^2=\frac{3+\sqrt{3}}{2}$.
\end{solution}
\end{question}

\begin{question}[points=4]
双曲线 $C:\frac{x^2}{a^2}-\frac{y^2}{b^2}=1$ 过点 $(2,\sqrt{3})$, 且离心率为 $2$, 则该双曲线的标准方程为（ ）
\begin{choices}
  \item $x^2-\frac{y^2}{3}=1$
  \item $\frac{x^2}{3}-y^2=1$
  \item $x^2-\frac{3y^2}{3}=1$
  \item $\frac{3x^2}{3}-y^2=1$
\end{choices}
\begin{solution}
\textbf{答案：}A

\textbf{分析：}分析可得 $b=\sqrt{3}a$, 再将点 $(2,\sqrt{3})$ 代入双曲线的方程, 求出 $a$ 的值, 即可得出双曲线的标准方程.

\textbf{详解：}$e=\frac{c}{a}=2$, 则 $c=2a$, $b=\sqrt{c^2-a^2}=\sqrt{3}a$, 则双曲线的方程为 $\frac{x^2}{a^2}-\frac{y^2}{3a^2}=1$, 将点 $(2,\sqrt{3})$ 代入可得 $\frac{4}{a^2}-\frac{3}{3a^2}=1$, 解得 $a^2=1$, 故 $b^2=3$, 因此双曲线的方程为 $x^2-\frac{y^2}{3}=1$.
\end{solution}
\end{question}

\begin{question}[points=4]
$\{a_n\}$ 和 $\{b_n\}$ 是两个等差数列, 其中 $\frac{a_k}{b_k}$ ($1\leq k\leq 5$) 为常值, $a_1=288$, $a_5=96$, $b_1=192$, 则 $b_3=$（ ）
\begin{choices}
  \item $64$
  \item $128$
  \item $256$
  \item $512$
\end{choices}
\begin{solution}
\textbf{答案：}B

\textbf{分析：}由已知条件求出 $b_5$ 的值, 利用等差中项的性质可求得 $b_3$ 的值.

\textbf{详解：}由已知条件可得 $\frac{a_1}{b_1}=\frac{a_5}{b_5}$, 则 $b_5=\frac{a_5 b_1}{a_1}=\frac{96\times192}{288}=64$, 因此 $b_3=\frac{b_1+b_5}{2}=\frac{192+64}{2}=128$.
\end{solution}
\end{question}

\begin{question}[points=4]
函数 $f(x)=\cos x-\cos 2x$, 试判断函数的奇偶性及最大值（ ）
\begin{choices}
  \item 奇函数, 最大值为 $2$
  \item 偶函数, 最大值为 $2$
  \item 奇函数, 最大值为 $\frac{9}{8}$
  \item 偶函数, 最大值为 $\frac{9}{8}$
\end{choices}
\begin{solution}
\textbf{答案：}D

\textbf{分析：}由函数奇偶性的定义结合三角函数的性质可判断奇偶性; 利用二倍角公式结合二次函数的性质可判断最大值.

\textbf{详解：}由题意, $f(-x)=\cos(-x)-\cos(-2x)=\cos x-\cos 2x=f(x)$, 所以该函数为偶函数, 又 $f(x)=\cos x-\cos 2x=-2\cos^2 x+\cos x+1=-2\left(\cos x-\frac{1}{4}\right)^2+\frac{9}{8}$, 所以当 $\cos x=\frac{1}{4}$ 时, $f(x)$ 取最大值 $\frac{9}{8}$.
\end{solution}
\end{question}

\begin{question}[points=4]
定义: 24小时内降水在平地上积水厚度（mm）来判断降雨程度. 其中小雨（$<10$ mm）, 中雨（$10$ mm $-25$ mm）, 大雨（$25$ mm $-50$ mm）, 暴雨（$50$ mm $-100$ mm）. 小明用一个圆锥形容器接了24小时的雨水, 如图, 则这天降雨属于哪个等级（ ）

\begin{center}
  \includegraphics[width=.38\linewidth]{question-8-1.jpg}
\end{center}
\begin{choices}
  \item 小雨
  \item 中雨
  \item 大雨
  \item 暴雨
\end{choices}
\begin{solution}
\textbf{答案：}B

\textbf{分析：}计算出圆锥体积, 除以圆面的面积即可得降雨量, 即可得解.

\textbf{详解：}由题意, 一个半径为 $\frac{200}{2}=100$ mm 的圆面内的降雨充满一个底面半径为 $\frac{200}{2}\times\frac{150}{300}=50$ mm, 高为 $150$ mm 的圆锥, 所以积水厚度 $d=\frac{\frac{1}{3}\pi\times 50^2 \times 150}{\pi\times 100^2}=12.5$ mm, 属于中雨.
\end{solution}
\end{question}

\begin{question}[points=4]
已知圆 $C:x^2+y^2=4$, 直线 $l:y=kx+m$, 当 $k$ 变化时, $l$ 截得圆 $C$ 弦长的最小值为 $2$, 则 $m=$（ ）
\begin{choices}
  \item $\pm 2$
  \item $\pm \sqrt{2}$
  \item $\pm \sqrt{3}$
  \item $\pm \sqrt{5}$
\end{choices}
\begin{solution}
\textbf{答案：}C

\textbf{分析：}先求得圆心到直线距离, 即可表示出弦长, 根据弦长最小值得出 $m$.

\textbf{详解：}由题可得圆心为 $(0,0)$, 半径为 $2$, 则圆心到直线的距离 $d=\frac{|m|}{\sqrt{k^2+1}}$, 则弦长为 $2\sqrt{4-\frac{m^2}{k^2+1}}$, 则当 $k=0$ 时, 弦长取得最小值为 $2\sqrt{4-m^2}=2$, 解得 $m=\pm\sqrt{3}$.
\end{solution}
\end{question}

\begin{question}[points=4]
数列 $\{a_n\}$ 是递增的整数数列, 且 $a_1\geq 3$, $a_1+a_2+\cdots+a_n=100$, 则 $n$ 的最大值为（ ）
\begin{choices}
  \item $9$
  \item $10$
  \item $11$
  \item $12$
\end{choices}
\begin{solution}
\textbf{答案：}C

\textbf{分析：}使数列首项、递增幅度均最小, 结合等差数列的通项及求和公式即可得解.

\textbf{详解：}若要使 $n$ 尽可能的大, 则 $a_1$, 递增幅度要尽可能小, 不妨设数列 $\{a_n\}$ 是首项为 $3$, 公差为 $1$ 的等差数列, 其前 $n$ 项和为 $S_n$, 则 $a_n=n+2$, $S_{11}=\frac{3+13}{2}\times 11=88<100$, $S_{12}=\frac{3+14}{2}\times 12=102>100$, 所以 $n$ 的最大值为 $11$.
\end{solution}
\end{question}

\section*{二、填空题}
本题共5小题, 每小题5分, 共25分.

\begin{question}[points=5]
$(x^3-\frac{1}{x})^4$ 展开式中常数项为.
\fillin[{$-4$}]
\begin{solution}
\textbf{答案：}$-4$

\textbf{详解：}$(x^3-\frac{1}{x})^4$ 的展开式的通项 $T_{r+1}=C_4^r (x^3)^{4-r}(-\frac{1}{x})^r=(-1)^r C_4^r x^{12-4r}$, 令 $r=3$ 得常数项为 $T_4=(-1)^3 C_4^3=-4$.
\end{solution}
\end{question}

\begin{question}[points=5]
已知抛物线 $C:y^2=4x$, 焦点为 $F$, 点 $M$ 为抛物线 $C$ 上的点, 且 $|FM|=6$, 则 $M$ 的横坐标是；作 $MN\perp x$ 轴于 $N$, 则 $S_{\triangle FMN}=$ .
\fillin[{$5$；$4\sqrt{5}$}]
\begin{solution}
\textbf{答案：}$5$；$4\sqrt{5}$

\textbf{分析：}根据焦半径公式可求 $M$ 的横坐标, 求出纵坐标后可求 $S_{\triangle FMN}$.

\textbf{详解：}因为抛物线的方程为 $y^2=4x$, 故 $p=2$ 且 $F(1,0)$. 因为 $|MF|=6$, $x_M+\frac{p}{2}=6$, 解得 $x_M=5$, 故 $y_M=\pm 2\sqrt{5}$, 所以 $S_{\triangle FMN}=\frac{1}{2}\times(5-1)\times 2\sqrt{5}=4\sqrt{5}$.
\end{solution}
\end{question}

\begin{question}[points=5]
$\vec{a}=(2,1)$, $\vec{b}=(2,-1)$, $\vec{c}=(0,1)$, 则 $(\vec{a}+\vec{b})\cdot\vec{c}=$ ； $\vec{a}\cdot\vec{b}=$ .
\fillin[{$0$；$3$}]
\begin{solution}
\textbf{答案：}$0$；$3$

\textbf{分析：}根据坐标求出 $\vec{a}+\vec{b}$, 再根据数量积的坐标运算直接计算即可.

\textbf{详解：}$\vec{a}+\vec{b}=(4,0)$, $\vec{a}\cdot\vec{c}=4\times0+0\times1=0$, $\vec{a}\cdot\vec{b}=2\times2+1\times(-1)=3$.
\end{solution}
\end{question}

\begin{question}[points=5]
若点 $P(\cos\theta,\sin\theta)$ 与点 $Q(\cos(\theta+\frac{\pi}{6}),\sin(\theta+\frac{\pi}{6}))$ 关于 $y$ 轴对称, 写出一个符合题意的 $\theta=$ .
\fillin[{$\frac{5\pi}{12}$（满足 $\theta=\frac{5\pi}{12}+k\pi, k\in\mathbb{Z}$ 即可）}]
\begin{solution}
\textbf{答案：}$\frac{5\pi}{12}$（满足 $\theta=\frac{5\pi}{12}+k\pi, k\in\mathbb{Z}$ 即可）

\textbf{分析：}根据 $P,Q$ 在单位圆上, 可得 $\theta,\theta+\frac{\pi}{6}$ 关于 $y$ 轴对称, 得出 $\theta+\frac{\pi}{6}+\theta=\pi+2k\pi, k\in\mathbb{Z}$ 求解.

\textbf{详解：}$\because P(\cos\theta,\sin\theta)$ 与 $Q(\cos(\theta+\frac{\pi}{6}),\sin(\theta+\frac{\pi}{6}))$ 关于 $y$ 轴对称, 即 $\theta,\theta+\frac{\pi}{6}$ 关于 $y$ 轴对称, $\theta+\frac{\pi}{6}+\theta=\pi+2k\pi, k\in\mathbb{Z}$, 则 $\theta=k\pi+\frac{5\pi}{12}, k\in\mathbb{Z}$, 当 $k=0$ 时, 可取 $\theta$ 的一个值为 $\frac{5\pi}{12}$.
\end{solution}
\end{question}

\begin{question}[points=5]
已知函数 $f(x)=|\lg x|-kx-2$, 给出下列四个结论：

①若 $k=0$, 则 $f(x)$ 有两个零点；

② $\exists k<0$, 使得 $f(x)$ 有一个零点；

③ $\exists k<0$, 使得 $f(x)$ 有三个零点；

④ $\exists k>0$, 使得 $f(x)$ 有三个零点.

以上正确结论的序号是.
\fillin[{$①②④$}]
\begin{solution}
\textbf{答案：}$①②④$

\textbf{分析：}由 $f(x)=0$ 可得出 $|\lg x|=kx+2$, 考查直线 $y=kx+2$ 与曲线 $g(x)=|\lg x|$ 的左、右支分别相切的情形, 利用方程思想以及数形结合可判断各选项的正误.

\textbf{详解：}对于①, 当 $k=0$ 时, 由 $f(x)=|\lg x|-2=0$, 可得 $x=\frac{1}{100}$ 或 $x=100$, ①正确；
对于②, 考查直线 $y=kx+2$ 与曲线 $y=-\lg x(0<x<1)$ 相切于点 $P(t,-\lg t)$, 对函数 $y=-\lg x$ 求导得 $y'=-\frac{1}{x\ln 10}$, 由题意可得 $\begin{cases} kt+2=-\lg t \\ k=-\frac{1}{t\ln 10} \end{cases}$, 解得 $\begin{cases} t=\frac{e}{100} \\ k=-\frac{100}{e}\lg e<0 \end{cases}$, 所以存在 $k=-\frac{100}{e}\lg e<0$, 使得 $f(x)$ 只有一个零点, ②正确；
对于③, 当直线 $y=kx+2$ 过点 $(1,0)$ 时, $k+2=0$, 解得 $k=-2$, 所以当 $-\frac{100}{e}\lg e<k<-2$ 时, 直线 $y=kx+2$ 与曲线 $y=-\lg x(0<x<1)$ 有两个交点, 若函数 $f(x)$ 有三个零点, 则直线 $y=kx+2$ 与曲线 $y=-\lg x(0<x<1)$ 有两个交点, 直线 $y=kx+2$ 与曲线 $y=\lg x(x>1)$ 有一个交点, 所以 $\begin{cases} -\frac{100}{e}\lg e<k<-2 \\ k+2>0 \end{cases}$, 此不等式无解, 因此不存在 $k<0$, 使得函数 $f(x)$ 有三个零点, ③错误；
对于④, 考查直线 $y=kx+2$ 与曲线 $y=\lg x(x>1)$ 相切于点 $P(t,\lg t)$, 对函数 $y=\lg x$ 求导得 $y'=\frac{1}{x\ln 10}$, 由题意可得 $\begin{cases} kt+2=\lg t \\ k=\frac{1}{t\ln 10} \end{cases}$, 解得 $\begin{cases} t=100e \\ k=\frac{\lg e}{100e} \end{cases}$, 所以当 $0<k<\frac{\lg e}{100e}$ 时, 函数 $f(x)$ 有三个零点, ④正确.
\end{solution}
\end{question}

\section*{三、解答题}
本题共6小题, 共85分. 解答应写出文字说明, 演算步骤或证明过程.

\begin{question}[points=13]
已知在 $\triangle ABC$ 中, $c=2b\cos B$, $C=\frac{2\pi}{3}$.

(1) 求 $B$ 的大小;

(2) 在下列三个条件中选择一个作为已知, 使 $\triangle ABC$ 存在且唯一确定, 并求出 $BC$ 边上的中线的长度.

① $c=\sqrt{2}b$; ② 周长为 $4+2\sqrt{3}$; ③ 面积为 $S_{\triangle ABC}=\frac{3\sqrt{3}}{4}$.
\begin{solution}
\textbf{答案：}(1) $B=\frac{\pi}{6}$; (2) 若选择②, 中线长为 $\sqrt{7}$; 若选择③, 中线长为 $\frac{\sqrt{21}}{2}$.

\textbf{分析：}(1) 由正弦定理化边为角即可求解; (2) 分别分析每个条件, 选择使三角形存在且唯一确定的, 利用余弦定理求解中线长.

\textbf{详解：}(1) 由 $c=2b\cos B$, 正弦定理得 $\sin C=2\sin B\cos B$, 即 $\sin 2B=\sin C=\sin\frac{2\pi}{3}=\frac{\sqrt{3}}{2}$, 又 $C=\frac{2\pi}{3}$, $B\in(0,\frac{\pi}{3})$, $2B\in(0,\frac{2\pi}{3})$, 所以 $2B=\frac{\pi}{3}$, 解得 $B=\frac{\pi}{6}$.
(2) 若选择①, 由正弦定理 $\frac{c}{b}=\frac{\sin C}{\sin B}=\frac{\sqrt{3}/2}{1/2}=\sqrt{3}$, 与 $c=\sqrt{2}b$ 矛盾, 故不存在.
若选择②, 由(1)得 $A=\pi-\frac{2\pi}{3}-\frac{\pi}{6}=\frac{\pi}{6}$, 则 $a=b$. 设外接圆半径为 $R$, 则 $a=b=2R\sin\frac{\pi}{6}=R$, $c=2R\sin\frac{2\pi}{3}=\sqrt{3}R$, 周长 $a+b+c=2R+\sqrt{3}R=4+2\sqrt{3}$, 解得 $R=2$, 则 $a=2$, $c=2\sqrt{3}$. 由余弦定理, $BC$ 边上的中线长 $m=\sqrt{b^2+(\frac{a}{2})^2-2\cdot b\cdot\frac{a}{2}\cdot\cos C}=\sqrt{(2)^2+1^2-2\times2\times1\times\cos\frac{2\pi}{3}}=\sqrt{4+1+2}=\sqrt{7}$.
若选择③, 由(1)得 $A=\frac{\pi}{6}$, 则 $a=b$. 面积 $S=\frac{1}{2}ab\sin C=\frac{1}{2}a^2\times\frac{\sqrt{3}}{2}=\frac{3\sqrt{3}}{4}$, 解得 $a=3$, 则 $b=3$, $c=2b\cos B=2\times3\times\frac{\sqrt{3}}{2}=3\sqrt{3}$. 由余弦定理, $BC$ 边上的中线长 $m=\sqrt{b^2+(\frac{a}{2})^2-2\cdot b\cdot\frac{a}{2}\cdot\cos C}=\sqrt{9+\frac{9}{4}-2\times3\times\frac{3}{2}\times\cos\frac{2\pi}{3}}=\sqrt{9+\frac{9}{4}+\frac{9}{2}}=\sqrt{\frac{36+9+18}{4}}=\sqrt{\frac{63}{4}}=\frac{3\sqrt{7}}{2}$.
\end{solution}
\end{question}

\begin{question}[points=14]
已知正方体 $ABCD-A_1B_1C_1D_1$, 点 $E$ 为 $A_1D_1$ 中点, 直线 $B_1C_1$ 交平面 $CDE$ 于点 $F$.

(1) 证明: 点 $F$ 为 $B_1C_1$ 的中点;

(2) 若点 $M$ 为棱 $A_1B_1$ 上一点, 且二面角 $M-CF-E$ 的余弦值为 $\frac{\sqrt{5}}{3}$, 求 $\frac{A_1M}{A_1B_1}$ 的值.
\begin{solution}
\textbf{答案：}(1) 证明见解析; (2) $\frac{A_1M}{A_1B_1}=\frac{1}{2}$.

\textbf{分析：}(1) 将平面 $CDE$ 扩展, 结合所得平面与直线 $B_1C_1$ 的交点即可证得; (2) 建立空间直角坐标系, 利用向量法求解.

\textbf{详解：}(1) 取 $B_1C_1$ 的中点 $F'$, 连接 $DE, EF', F'C$, 由正方体性质, $E,F'$ 为中点, 故 $EF'\parallel CD$, 从而 $E,F',C,D$ 四点共面, 即平面 $CDE$ 即平面 $CDEF'$, 所以直线 $B_1C_1$ 交平面 $CDE$ 于点 $F'$, 由直线与平面相交唯一交点, 故点 $F$ 与点 $F'$ 重合, 即点 $F$ 为 $B_1C_1$ 中点.
(2) 以点 $D$ 为原点, $DA, DC, DD_1$ 方向为 $x,y,z$ 轴建立空间直角坐标系, 设正方体棱长为 $2$, 设 $\frac{A_1M}{A_1B_1}=\lambda(0\leq\lambda\leq1)$, 则 $M(2,2\lambda,2), C(0,2,0), F(1,2,2), E(1,0,2)$. 则 $\overrightarrow{MC}=(-2,2-2\lambda,-2)$, $\overrightarrow{CF}=(1,0,2)$, $\overrightarrow{FE}=(0,-2,0)$. 设平面 $MCF$ 的法向量 $\vec{m}=(x_1,y_1,z_1)$, 由 $\begin{cases} \vec{m}\cdot\overrightarrow{MC}=-2x_1+(2-2\lambda)y_1-2z_1=0 \\ \vec{m}\cdot\overrightarrow{CF}=x_1+2z_1=0 \end{cases}$, 令 $z_1=-1$, 得 $\vec{m}=(2,\frac{1}{1-\lambda},-1)$. 设平面 $CFE$ 的法向量 $\vec{n}=(x_2,y_2,z_2)$, 由 $\begin{cases} \vec{n}\cdot\overrightarrow{FE}=-2y_2=0 \\ \vec{n}\cdot\overrightarrow{CF}=x_2+2z_2=0 \end{cases}$, 令 $z_2=-1$, 得 $\vec{n}=(2,0,-1)$. 则 $\cos\langle\vec{m},\vec{n}\rangle=\frac{\vec{m}\cdot\vec{n}}{|\vec{m}||\vec{n}|}=\frac{5}{\sqrt{5+(\frac{1}{1-\lambda})^2}\cdot\sqrt{5}}=\frac{\sqrt{5}}{3}$, 解得 $(\lambda-1)^2=\frac{1}{4}$, 故 $\lambda=\frac{1}{2}$ (舍去 $\frac{3}{2}$). 所以 $\frac{A_1M}{A_1B_1}=\frac{1}{2}$.
\end{solution}
\end{question}

\begin{question}[points=14]
为加快新冠肺炎检测效率, 某检测机构采取“$k$ 合 $1$ 检测法”, 即将 $k$ 个人的拭子样本合并检测, 若为阴性, 则可以确定所有样本都是阴性的; 若为阳性, 则还需要对本组的每个人再做检测. 现有 $100$ 人, 已知其中 $2$ 人感染病毒.

(1) ①若采用“$10$ 合 $1$ 检测法”, 且两名患者在同一组, 求总检测次数;

②已知 $10$ 人分成一组, 分 $10$ 组, 两名感染患者在同一组的概率为 $\frac{1}{11}$, 定义随机变量 $X$ 为总检测次数, 求检测次数 $X$ 的分布列和数学期望 $E(X)$;

(2) 若采用“$5$ 合 $1$ 检测法”, 检测次数 $Y$ 的期望为 $E(Y)$, 试比较 $E(X)$ 和 $E(Y)$ 的大小(直接写出结果).
\begin{solution}
\textbf{答案：}(1) ① $20$ 次; ②分布列见解析, 期望为 $\frac{320}{11}$; (2) 见解析.

\textbf{分析：}(1) ①由题设条件还原情境, 即可得解; ②求出 $X$ 的取值情况及概率, 进而可得分布列和期望; (2) 求出 $E(Y)$, 分类比较.

\textbf{详解：}(1) ①对每组进行检测, 需要 $10$ 次; 再对结果为阳性的组每个人进行检测, 需要 $10$ 次; 所以总检测次数为 $20$ 次.
②由题意, $X$ 可以取 $20,30$, $P(X=20)=\frac{1}{11}$, $P(X=30)=1-\frac{1}{11}=\frac{10}{11}$, 分布列为:
| $X$ | $20$ | $30$ |
|-----|-----|-----|
| $P$ | $\frac{1}{11}$ | $\frac{10}{11}$ |
所以 $E(X)=20\times\frac{1}{11}+30\times\frac{10}{11}=\frac{320}{11}$.
(2) 由题意, $Y$ 可以取 $25,30$, 设两名感染者在同一组的概率为 $p$, $P(Y=25)=p$, $P(Y=30)=1-p$, 则 $E(Y)=25p+30(1-p)=30-5p$. 若 $p=\frac{2}{11}$, $E(X)=E(Y)$; 若 $p>\frac{2}{11}$, $E(X)>E(Y)$; 若 $p<\frac{2}{11}$, $E(X)<E(Y)$.
\end{solution}
\end{question}

\begin{question}[points=15]
已知函数 $f(x)=\frac{3-2x}{x^2+a}$.

(1) 若 $a=0$, 求 $y=f(x)$ 在 $(1,f(1))$ 处的切线方程;

(2) 若函数 $f(x)$ 在 $x=-1$ 处取得极值, 求 $f(x)$ 的单调区间, 以及最大值和最小值.
\begin{solution}
\textbf{答案：}(1) $4x+y-5=0$; (2) 增区间为 $(-\infty,-1)$ 和 $(4,+\infty)$, 减区间为 $(-1,4)$, 最大值为 $1$, 最小值为 $-\frac{1}{4}$.

\textbf{分析：}(1) 求出 $f(1),f'(1)$ 的值, 利用点斜式写出切线方程; (2) 由 $f'(-1)=0$ 求 $a$, 然后利用导数分析单调性与极值.

\textbf{详解：}(1) 当 $a=0$ 时, $f(x)=\frac{3-2x}{x^2}$, $f'(x)=\frac{2(x-3)}{x^3}$, $f(1)=1$, $f'(1)=-4$, 切线方程为 $y-1=-4(x-1)$, 即 $4x+y-5=0$.
(2) $f(x)=\frac{3-2x}{x^2+a}$, $f'(x)=\frac{-2(x^2+a)-2x(3-2x)}{(x^2+a)^2}=\frac{2(x^2-3x-a)}{(x^2+a)^2}$, 由 $f'(-1)=0$ 得 $\frac{2(1+3-a)}{(1+a)^2}=0$, 解得 $a=4$. 故 $f(x)=\frac{3-2x}{x^2+4}$, $f'(x)=\frac{2(x+1)(x-4)}{(x^2+4)^2}$. 列表:
| $x$ | $(-\infty,-1)$ | $-1$ | $(-1,4)$ | $4$ | $(4,+\infty)$ |
|-----|----------------|------|----------|-----|---------------|
| $f'(x)$ | $+$ | $0$ | $-$ | $0$ | $+$ |
| $f(x)$ | 增 | 极大值 | 减 | 极小值 | 增 |
所以增区间为 $(-\infty,-1)$ 和 $(4,+\infty)$, 减区间为 $(-1,4)$. 当 $x<\frac{3}{2}$ 时 $f(x)>0$, $x>\frac{3}{2}$ 时 $f(x)<0$, 所以 $f(x)_{\max}=f(-1)=1$, $f(x)_{\min}=f(4)=-\frac{1}{4}$.
\end{solution}
\end{question}

\begin{question}[points=14]
已知椭圆 $E:\frac{x^2}{a^2}+\frac{y^2}{b^2}=1(a>b>0)$ 过点 $A(0,-2)$, 以四个顶点围成的四边形面积为 $4\sqrt{5}$.

(1) 求椭圆 $E$ 的标准方程;

(2) 过点 $P(0,-3)$ 的直线 $l$ 斜率为 $k$, 交椭圆 $E$ 于不同的两点 $B,C$, 直线 $AB,AC$ 交 $y=-3$ 于点 $M,N$, 若 $|PM|+|PN|\leq15$, 求 $k$ 的取值范围.
\begin{solution}
\textbf{答案：}(1) $\frac{x^2}{5}+\frac{y^2}{4}=1$; (2) $[-3,-1)\cup(1,3]$.

\textbf{分析：}(1) 根据椭圆过点及四边形面积可求 $a,b$; (2) 设 $B(x_1,y_1),C(x_2,y_2)$, 求出 $M,N$ 的横坐标, 利用弦长公式和韦达定理化简 $|PM|+|PN|$, 结合判别式可得范围.

\textbf{详解：}(1) 因为椭圆过 $A(0,-2)$, 故 $b=2$, 四个顶点围成的四边形面积为 $\frac{1}{2}\times2a\times2b=4\sqrt{5}$, 即 $a=\sqrt{5}$, 所以椭圆方程为 $\frac{x^2}{5}+\frac{y^2}{4}=1$.
(2) 设 $B(x_1,y_1),C(x_2,y_2)$, 直线 $BC:y=kx-3$, 联立 $\begin{cases} y=kx-3 \\ 4x^2+5y^2=20 \end{cases}$ 得 $(4+5k^2)x^2-30kx+25=0$, $\Delta=900k^2-100(4+5k^2)>0$, 解得 $k<-1$ 或 $k>1$, 且 $x_1+x_2=\frac{30k}{4+5k^2}$, $x_1x_2=\frac{25}{4+5k^2}>0$. 直线 $AB:y=\frac{y_1+2}{x_1}x-2$, 令 $y=-3$ 得 $x_M=-\frac{x_1}{y_1+2}$, 同理 $x_N=-\frac{x_2}{y_2+2}$. 所以 $|PM|+|PN|=|x_M|+|x_N|=\left|\frac{x_1}{y_1+2}\right|+\left|\frac{x_2}{y_2+2}\right|=\frac{|x_1|}{|kx_1-1|}+\frac{|x_2|}{|kx_2-1|}$. 由于 $x_1x_2>0$, 且 $k<-1$ 或 $k>1$ 时 $kx_i-1$ 同号, 故 $|PM|+|PN|=\left|\frac{x_1}{kx_1-1}+\frac{x_2}{kx_2-1}\right|=\left|\frac{2kx_1x_2-(x_1+x_2)}{k^2x_1x_2-k(x_1+x_2)+1}\right|=\left|\frac{2k\cdot\frac{25}{4+5k^2}-\frac{30k}{4+5k^2}}{k^2\cdot\frac{25}{4+5k^2}-k\cdot\frac{30k}{4+5k^2}+1}\right|=\left|\frac{50k-30k}{25k^2-30k^2+4+5k^2}\right|=\left|\frac{20k}{4}\right|=5|k|$. 由 $|PM|+|PN|\leq15$ 得 $5|k|\leq15$, 即 $|k|\leq3$, 结合 $k<-1$ 或 $k>1$, 得 $k$ 的取值范围为 $[-3,-1)\cup(1,3]$.
\end{solution}
\end{question}

\begin{question}[points=15]
定义 $R_p$ 数列 $\{a_n\}$: 对实数 $p$, 满足: ① $a_1+p\geq0$, $a_2+p=0$; ② $\forall n\in\mathbb{N}^*$, $a_{4n-1}<a_{4n}$; ③ $a_{m+n}\in\{a_m+a_n+p, a_m+a_n+p+1\}$, $m,n\in\mathbb{N}^*$.

(1) 对于前 $4$ 项 $2,-2,0,1$ 的数列, 可以是 $R_2$ 数列吗? 说明理由;

(2) 若 $\{a_n\}$ 是 $R_0$ 数列, 求 $a_5$ 的值;

(3) 是否存在 $p$, 使得存在 $R_p$ 数列 $\{a_n\}$, 对 $\forall n\in\mathbb{N}^*$, $S_n\geq S_{10}$? 若存在, 求出所有这样的 $p$; 若不存在, 说明理由.
\begin{solution}
\textbf{答案：}(1) 不可以, 理由见解析; (2) $a_5=1$; (3) 存在, $p=2$.

\textbf{分析：}(1) 利用性质③验证; (2) 由性质推导前四项, 再用数学归纳法; (3) 构造 $b_n=a_n+p$, 利用(2)结论.

\textbf{详解：}(1) 由性质③, $a_3\in\{a_1+a_2+2, a_1+a_2+2+1\}=\{2,3\}$, 但 $a_3=0$ 不在其中, 矛盾, 故不可以是 $R_2$ 数列.
(2) 由 $R_0$ 性质, $a_1\geq0$, $a_2=0$. 由性质③, $a_3\in\{a_1, a_1+1\}$, $a_4\in\{0,1\}$. 若 $a_4=0$, 由性质② $a_3<a_4=0$, 则 $a_1<0$ 或 $a_1+1<0$, 矛盾; 若 $a_4=1$, $a_3=a_1+1$, 由 $a_3<1$ 得 $a_1<0$, 矛盾; 故 $a_4=1$, $a_3=a_1$. 再由 $a_4=a_1+a_3$ 或 $a_4=a_1+a_3+1$, 得 $a_1=\frac{1}{2}$ 或 $0$. 若 $a_1=\frac{1}{2}$, 则 $a_2=a_{1+1}\in\{2a_1,2a_1+1\}=\{1,2\}$, 与 $a_2=0$ 矛盾. 故 $a_1=0$, 前四项为 $0,0,0,1$. 用数学归纳法证明 $a_{4n+i}=n$ ($i=1,2,3$), $a_{4n+4}=n+1$, $n\in\mathbb{N}$. 当 $n=0$ 时成立. 假设 $n\leq k$ 时成立, 则 $n=k+1$ 时, 可推出 $a_{4k+5}=k+1$, $a_{4k+6}=k+1$, $a_{4k+8}=k+2$, $a_{4k+7}=k+1$. 故 $a_5=a_{4\times1+1}=1$.
(3) 令 $b_n=a_n+p$, 则 $b_n$ 满足 $R_0$ 性质, 由(2)知 $b_{4n+i}=n$ ($i=1,2,3$), $b_{4n+4}=n+1$. 故 $a_n=b_n-p$. 要求 $S_n\geq S_{10}$ 恒成立, 即 $S_n-S_{10}\geq0$, 考察 $S_{11}-S_{10}=a_{11}=b_{11}-p=2-p\geq0$, $S_9-S_{10}=-a_{10}=-(b_{10}-p)=-(2-p)\geq0$, 故 $p=2$. 此时 $a_1,a_2,\ldots,a_{10}\leq0$, $a_j\geq0$ ($j\geq11$), 满足题意. 所以存在 $p=2$.
\end{solution}
\end{question}

